\documentclass[../main.tex]{subfiles}

\begin{document}
\subsection{Teorema de Cauchy homotópico}
\classdate{2026-04-27}

\begin{definition}
    Sean \(\gamma_{0},\gamma_{1} \colon [a, b] \to \mathbb{C}\) dos curvas
    (continuas) en
    una región \(U\) de \(\mathbb{C}\). Una homotopía entre \(\gamma_{0}\) y
    \(\gamma_{1}\) en \(U\) es una transformación continua \(H \colon [a, b] \mul [0, 1]
    \to \mathbb{C}\) con las siguientes propiedades

    \begin{enumerate}
        \item \(H(t, 0) = \gamma_{0}(t),\medspace \forall t\in [a, b]\).
        \item \(H(t, 1) = \gamma_{1}(t),\medspace \forall t\in [a, b]\).
    \end{enumerate}

    Cuando existe tal homotopía, decimos que las curvas \(\gamma_{0},\gamma_{1}\) son
    homotópicas.
\end{definition}

En este contexto, definimos \(\gamma_{s}(t) = H(t, s)\), con \(s\in [0, 1]\) fija.
Podemos pensar a la familia de curvas \(\{\gamma_{s}\}\) como una deformación continua
de \(\gamma_{0}\) en \(\gamma_{1}\).

\begin{remark}
    La definición anterior requiere de varios ajustes para que nos sea realmente útil.
    Puesto que en nuestro contexto estamos integrando usualmente funciones analítica
    sobre curvas cerradas suaves por pedazos, por lo que tendremos los siguientes
    requisitos adicionales:

    \begin{itemize}
        \item Si \(\gamma_{0}\) y \(\gamma_{1}\) son curvas cerradas, es decir,

            \begin{equation*}
                \gamma_{0}(a) = \gamma_{0}(b)\quad \text{y}\quad
                \gamma_{1}(a) = \gamma_{1}(b)
            \end{equation*}

            pediremos que las curvas \(\gamma_{s}, \forall s\in [0, 1]\)
            también lo sean.

        \item Si \(\gamma_{0}\) y \(\gamma_{1}\) son curvas suaves por
            pedazos, queremos que
            \(\gamma_{s}, \forall s\in [0, 1]\) sean curvas cerradas por pedazos.
    \end{itemize}
\end{remark}

\begin{definition}
    Decimos que \(U \subseteq \mathbb{C}\) es simplemente conexa si y solo si toda curva
    cerrada \(\gamma \subseteq U\) es homotópica a algún \(z_{0} \in U\).
\end{definition}

\begin{theorem}[de Cauchy para homotopías]\label{thm:Cauchy-Homotopies}
    Sea \(U \subseteq \mathbb{C}\) una región, \(f \colon U \to \mathbb{C}\) una función
    analítica y \(\gamma_{0}, \gamma_{1} \colon [a, b] \to \mathbb{C}\) dos curvas
    cerradas suaves por pedazos. Si \(\gamma_{0}, \gamma_{1}\) son homotópicas, entonces

    \begin{equation*}
        \int_{\gamma_{0}}^{}f(z)\odif{z} = \int_{\gamma_{1}}^{}f(z)\odif{z}.
    \end{equation*}
\end{theorem}

Andas de dar la prueba, enunciaremos un resultado

\begin{corollary}
    Sea \(f \colon U \to \mathbb{C}\) analítica, \(U\) simplemente
    conexa\footnote{Si toda
    curva cerrada en \(U\) es homotópica.} y \(\gamma \colon [a, b] \to U\)
    cerrada suave
    por pedazos. Entonces,

    \begin{equation*}
        \int_{\gamma}^{}f(z)\odif{z} = 0.
    \end{equation*}
\end{corollary}

\begin{proof}[\zcref{thm:Cauchy-Homotopies}]
    Sea \(H \colon [a, b] \mul [0, 1] \to U\) una homotopía entre \(\gamma_{0}\) y
    \(\gamma_{1}\) y, para simplificar, asumiremos que \(H\) es de clase \(C^{2}\).
    Definimos \(h \colon [0, 1] \to \mathbb{C}\) como

    \begin{equation*}
        h(s) = \int_{\gamma_{0}}^{}f(z)\odif{z} = \int_{a}^{b}f(H(t, s))\pdv{H(s,
        t)}{t}\odif{t}
    \end{equation*}

    con \(\gamma_{s}(t) = H(t, s)\).

    Queremos demostrar que \(h\) es constante en \([0, 1]\). Para esto, mostraremos que
    su derivada en idénticamente igual a 0:

    \begin{align*}
        \odv{h}{s} &= \odv{}{s}\int_{a}^{b}f(H(t, s))\pdv{H(t, s)}{t}\odif{t} =
        \int_{a}^{b}\pdv*{\biggl(f(H(t, s))\pdv{H(t, s)}{t}\biggr)}{s}\odif{t},\\
        &= \int_{a}^{b}\Biggl[f^{\prime}(H(t, s))\pdv{H(t, s)}{s}\pdv{H(t, s)}{t} +
        f(H(t, s))\pdv{H(t, s)}{s, t}\Biggr]\odif{t},\\
        &= \int_{a}^{b}\Biggl[f^{\prime}(H(t, s))\pdv{H(t, s)}{t}\pdv{H(t, s)}{s} +
        f(H(t, s))\pdv{H(t, s)}{t,s}\Biggr]\odif{t},\\
        &= \Biggl[f(H(t, s))\pdv{H(t, s)}{s}\Biggr]_{t = a}^{t = b} = 0.
    \end{align*}
\end{proof}

\classdate{2026-04-28}

Enunciamos los siguientes corolarios de la versión homotópica del teorema de Cauchy.

\begin{corollary}
    Sean \(U \subseteq \mathbb{C}\) una región y \(f \colon \mathbb{C} \to \mathbb{C}\)
    analítica y \(\gamma \colon [a, b] \to U\) una curva cerrada suave por pedazos
    homotópica a una curva constante en \(U\), entonces

    \begin{equation*}
        \int_{\gamma}^{}f(z)\odif{z} = 0.
    \end{equation*}
\end{corollary}

\begin{corollary}
    Sean \(U \subseteq \mathbb{C}\) una región simplemente conexa\footnote{Toda curva es
    homotópica a un punto.} y \(f \colon  U  \subseteq \mathbb{C} \to \mathbb{C}\)
    analítica. Entonces, \( \forall \gamma\colon [a, b] \to U\) cerrada suave
    por pedazos,

    \begin{equation}
        \int_{\gamma}^{}f(z)\odif{z} = 0.
    \end{equation}
\end{corollary}

\begin{theorem}[Fórmula Integral de Cauchy]
    Sea \(U \subseteq \mathbb{C}\) una región y \(f \colon U \subseteq \mathbb{C} \to
    \mathbb{C}\) una función analítica. Sea \(\gamma \colon [a, b] \to U\) una curva
    cerrada suave por pedazos y sea \(z_{0}\in  U\setminus \gamma([a, b])\). Si
    \(\gamma\) es homotópica a una constante en \(U\), entonces

    \begin{equation*}
        n(\gamma, z_{0})f(z_{0}) = \dfrac{1}{2\pi i} \int_{\gamma}^{}\dfrac{f(z)}{z -
        z_{0}}\odif{z}.
    \end{equation*}
\end{theorem}

\begin{proof}
    Sea \(g \colon U \subseteq \mathbb{C} \to \mathbb{C}\) una función
    auxiliar dada como

    \begin{equation*}
        g(z) =
        \begin{dcases}
            \dfrac{f(z) - f(z_{0})}{z - z_{0}}, &\text{si}\; z \neq z_{0},\\
            f^{\prime}(z_{0}), &\text{si}\; z = z_{0}
        \end{dcases}
    \end{equation*}

    ¿Y qué podemos decir de la función?

    \(g\) es continua en \(z_{0}\), pues

    \begin{equation*}
        \lim_{z \to z_{0}} g(z) = \lim_{z \to z_{0}} \dfrac{f(z) -
        f(z_{0})}{z - z_{0}} =
        f^{\prime}(z_{0}).
    \end{equation*}

    Y tiene una singularidad removible en \(z_{0}\),

    \begin{equation*}
        \lim_{z \to z_{0}} \bigl(z - z_{0}\bigr)\dfrac{f(z) - f(z_{0})}{z - z_{0}} = 0.
    \end{equation*}

    Y por lo tanto,  la función admite una extensión analítica, que por
    continuidad, debe
    coincidir con \(g\). Por lo que \(g\) es analítica.

    Entonces por el teorema de Cauchy,

    \begin{equation*}
        \int_{\gamma}^{}g(z)\odif{z} = 0.
    \end{equation*}

    Como \(z\notin \gamma([a, b])\), entonces \(\forall z\in \gamma([a, b])\),

    \begin{equation*}
        \int_{\gamma}^{}\dfrac{f(z)}{z - z_{0}}\odif{z} =
        f(z_{0})\underbrace{\int_{\gamma}^{}\dfrac{1}{z - z_{0}}\odif{z}}_{=
            2\pi i n(\gamma,
        z_{0})}.
    \end{equation*}
\end{proof}

\begin{corollary}
    Si \(U\) es una región simplemente conexa y \(F \colon U \subseteq \mathbb{C} \to
    \mathbb{C}\) ese analítica, tal que \(F^{\prime} = f\) en \(U\).
\end{corollary}

\begin{proof}
    Sea \(\gamma \colon [a, b] \to U\) una curva cerrada suave por pedazos. Como
    \(\gamma\) es homotópica a una constante en \(U\), pues \(U\) es simplemente conexa;
    por el teorema de Cauchy,

    \begin{equation*}
        \int_{\gamma}^{}f(z)\odif{z} = 0.
    \end{equation*}

    Y por el teorema e la primitiva, \(f\) admite una primitiva \(F\) en \(U\).
\end{proof}

Ahora, desviamos un poco nuestra atención, para hablar un poco acerca de las funciones
armónicas.

Nuestro objetivo en esta parte es dar un principio del módulo máximo para funciones
armónicas; es decir, si \(U \subseteq \mathbb{C}\) una región y \(u \colon U \to
\mathbb{R}\) es una función armónica que alcanza su máximo en \(z_{0}\in  U\), entonces
\(u\) es constante.
La motivación detrás de esto se debe a que la parte real de una función analítica es
armónica. Si \(f = u + iv\), pues es analítica, satisface Cauchy-Riemann

\begin{equation*}
    \pdv{u}{x} = \pdv{v}{y}\quad \text{y}\quad \pdv{u}{y} = - \pdv{v}{x}.
\end{equation*}

Entonces,

\begin{equation*}
    \pdv{u}{x,x} = \pdv{v}{x,y}\quad \text{y}\quad \pdv{u}{y,y} = - \pdv{v}{y,x} = -
    \pdv{v}{x,y}.
\end{equation*}

Sumándolas,

\begin{equation*}
    \Delta{u} = \pdv{u}{x,x} + \pdv{u}{y,y} = 0.
\end{equation*}

Una pregunta que nos podemos hacer es: ¿se vale el recíproco? En caso de que sí,  para
\(u \colon U \to \mathbb{C}\) armónica tendríamos \(f \colon U \subseteq \mathbb{C} \to
\mathbb{C}\) analítica, tal que \(\Re(f) = u\). Entonces, \(\abs{\mathrm{e}^{f(z)}} =
\mathrm{e}^{\Re(f(z))} = e^{u(z)}\), y como \(\mathrm{e}^{x}\) es creciente,
\(\abs{\mathrm{e}^{f(z)}}\) alcanza su máximo en \(z_{0}\).
\end{document}
